\documentclass{beamer}

\usepackage{amsmath,amsfonts,amssymb}
\usepackage{physics}
\usepackage{bm}
\usepackage{quantikz}
\usepackage{listings}

\title{Quantum Fourier Transform and Quantum Phase Estimation}
\author{Morten Hjorth-Jensen}
\date{Additional notes spring 2026}

\lstset{
language=Python,
basicstyle=\ttfamily\small,
keywordstyle=\color{blue},
commentstyle=\color{green!60!black}
}

\begin{document}

\frame{\titlepage}

%================================================
\section{Overview}

\begin{frame}{Overview}

Lecture 1

\begin{itemize}
\item Fourier transforms in physics
\item Mathematical definition of QFT
\item Derivation and structure
\end{itemize}

Lecture 2

\begin{itemize}
\item QFT circuits
\item Approximate QFT
\item Complexity
\end{itemize}

Lecture 3

\begin{itemize}
\item Quantum Phase Estimation
\item Error analysis
\item Iterative phase estimation
\end{itemize}

Lecture 4

\begin{itemize}
\item Hamiltonian eigenvalue estimation
\item Quantum chemistry
\item Numerical simulations
\end{itemize}

\end{frame}

%================================================
\section{Lecture 1: Mathematical Foundations}

\begin{frame}{Fourier Analysis in Physics}

Fourier transforms appear in

\begin{itemize}
\item wave mechanics
\item signal processing
\item spectral decomposition
\item lattice models
\end{itemize}

Example

\[
\psi(x)=\int dk \; \tilde{\psi}(k)e^{ikx}
\]

\end{frame}

%------------------------------------------------

\begin{frame}{Discrete Fourier Transform}

For vector

\[
x=(x_0,\dots,x_{N-1})
\]

\[
y_k =
\frac{1}{\sqrt N}
\sum_{j=0}^{N-1}
x_j e^{2\pi i jk/N}
\]

\end{frame}

%------------------------------------------------

\begin{frame}{Fourier Matrix}

\[
F_{jk}=
\frac{1}{\sqrt N}
e^{2\pi i jk/N}
\]

Properties

\begin{itemize}
\item unitary
\item periodic
\item diagonalizes cyclic shifts
\end{itemize}

\end{frame}

%------------------------------------------------

\begin{frame}{Quantum States}

Quantum register

\[
|\psi\rangle =
\sum_x \alpha_x |x\rangle
\]

Dimension

\[
N=2^n
\]

\end{frame}

%------------------------------------------------

\begin{frame}{Definition of QFT}

\[
QFT|x\rangle =
\frac{1}{\sqrt N}
\sum_{k=0}^{N-1}
e^{2\pi i xk/N}|k\rangle
\]

\end{frame}

%------------------------------------------------

\begin{frame}{Linearity}

For arbitrary state

\[
|\psi\rangle =
\sum_x \alpha_x |x\rangle
\]

\[
QFT|\psi\rangle =
\sum_k \tilde{\alpha}_k |k\rangle
\]

\end{frame}

%------------------------------------------------

\begin{frame}{Binary Representation}

\[
x=x_{n-1}2^{n-1}+\dots+x_0
\]

\[
x/2^n=0.x_{n-1}x_{n-2}\dots x_0
\]

\end{frame}

%------------------------------------------------

\begin{frame}{Phase Expansion}

\[
e^{2\pi i xk/2^n}
\]

can be decomposed into binary fractions.

This leads to factorization of the QFT.

\end{frame}

%------------------------------------------------

\begin{frame}{Factorized Form}

\[
QFT|x\rangle=
\frac{1}{2^{n/2}}
\bigotimes_{j=0}^{n-1}
(|0\rangle+
e^{2\pi i(0.x_j\dots x_0)}|1\rangle)
\]

\end{frame}

%------------------------------------------------

\begin{frame}{Interpretation}

Each qubit accumulates phases from less significant bits.

Controlled rotations naturally appear.

\end{frame}

%================================================
\section{Lecture 2: QFT Circuits}

\begin{frame}{Phase Rotation Gates}

Define

\[
R_k=
\begin{pmatrix}
1&0\\
0&e^{2\pi i/2^k}
\end{pmatrix}
\]

\end{frame}

%------------------------------------------------

\begin{frame}{General QFT Circuit}

\begin{center}
\begin{quantikz}
\lstick{q0}&\gate{H}&\ctrl{1}&\ctrl{2}&\swap{2}&\qw\\
\lstick{q1}&\qw&\gate{R2}&\ctrl{1}&\qw&\qw\\
\lstick{q2}&\qw&\qw&\gate{R3}&\gate{H}&\swap{-2}
\end{quantikz}
\end{center}

\end{frame}


%------------------------------------------------

\begin{frame}{Two-Qubit QFT}

\begin{center}
\begin{quantikz}
\lstick{}&\gate{H}&\ctrl{1}&\swap{1}\\
\lstick{}&\gate{R2}&\gate{H}&\swap{-1}
\end{quantikz}
\end{center}

\end{frame}

%------------------------------------------------

\begin{frame}{Three-Qubit QFT}

\begin{center}
\begin{quantikz}
\lstick{}&\gate{H}&\ctrl{1}&\ctrl{2}&\swap{2}\\
\lstick{}&\gate{R2}&\gate{H}&\ctrl{1}&\qw\\
\lstick{}&\gate{R3}&\gate{R2}&\gate{H}&\swap{-2}
\end{quantikz}
\end{center}

\end{frame}

%------------------------------------------------

\begin{frame}{Gate Complexity}

Number of gates

\[
\frac{n(n+1)}{2}
\]

Scaling

\[
O(n^2)
\]

\end{frame}

%------------------------------------------------

\begin{frame}{Approximate QFT}

Small rotations contribute little.

Truncate rotations

\[
R_k \quad k>m
\]

\end{frame}

%------------------------------------------------

\begin{frame}{Complexity Reduction}

Approximate QFT

\[
O(n\log n)
\]

\end{frame}

%================================================
\section{Lecture 3: Quantum Phase Estimation}

\begin{frame}{Eigenvalue Problem}

Suppose

\[
U|\psi\rangle=e^{2\pi i\phi}|\psi\rangle
\]

Goal

\[
\phi
\]

\end{frame}

%------------------------------------------------

\begin{frame}{Registers}

Two registers

\begin{itemize}
\item phase register
\item eigenstate register
\end{itemize}

\end{frame}

%------------------------------------------------

\begin{frame}{Initial State}

\[
|0\rangle^{\otimes t}|\psi\rangle
\]

\end{frame}

%------------------------------------------------

\begin{frame}{Superposition}

Apply Hadamards

\[
\frac{1}{2^{t/2}}
\sum_k |k\rangle |\psi\rangle
\]

\end{frame}

%------------------------------------------------

\begin{frame}{Controlled Powers}

Controlled

\[
U^{2^k}
\]

\end{frame}

%------------------------------------------------

\begin{frame}{State Before QFT}

\[
\frac{1}{2^{t/2}}
\sum_k e^{2\pi i k\phi}|k\rangle|\psi\rangle
\]

\end{frame}

%------------------------------------------------

\begin{frame}{Inverse QFT}

Inverse QFT converts phase to binary digits.

\end{frame}

%------------------------------------------------

\begin{frame}{Measurement}

Measurement yields

\[
\phi \approx 0.\phi_1\phi_2\dots
\]

\end{frame}

%================================================
\section{Error Analysis}

\begin{frame}{Finite Precision}

If

\[
\phi\neq k/2^t
\]

measurement yields nearest integer.

\end{frame}

%------------------------------------------------

\begin{frame}{Probability Distribution}

\[
P(y)=
\frac{1}{2^{2t}}
\left|
\sum_k e^{2\pi i k(\phi-y/2^t)}
\right|^2
\]

\end{frame}

%------------------------------------------------

\begin{frame}{Success Bound}

Probability

\[
P \ge 4/\pi^2
\]

\end{frame}

%================================================
\section{Iterative Phase Estimation}

\begin{frame}{Motivation}

Standard QPE requires many qubits.

Iterative version uses one ancilla.

\end{frame}

%------------------------------------------------

\begin{frame}{Iterative Circuit}

\begin{center}
\begin{quantikz}
\lstick{}&\gate{H}&\ctrl{1}&\meter{}\\
\lstick{psi}&\qw&\gate{U2k}&\qw
\end{quantikz}
\end{center}

\end{frame}

%================================================
\section{Hamiltonian Eigenvalue Estimation}

\begin{frame}{Hamiltonian Simulation}

If

\[
U=e^{-iHt}
\]

then

\[
U|\psi_k\rangle=e^{-iE_k t}|\psi_k\rangle
\]

\end{frame}

%------------------------------------------------

\begin{frame}{Energy Extraction}

\[
E_k =
\frac{2\pi\phi}{t}
\]

\end{frame}

%------------------------------------------------

\begin{frame}{Applications}

\begin{itemize}
\item quantum chemistry
\item materials science
\item nuclear physics
\end{itemize}

\end{frame}

%================================================
\section{Connection to VQE}

\begin{frame}{Comparison}

Phase estimation

\begin{itemize}
\item exact
\item deep circuits
\end{itemize}

VQE

\begin{itemize}
\item shallow circuits
\item classical optimization
\end{itemize}

\end{frame}

%================================================
\section{Python Simulation}

\begin{frame}[fragile]{QFT Matrix}

\begin{lstlisting}
import numpy as np

def qft_matrix(n):

    N=2**n
    omega=np.exp(2j*np.pi/N)

    F=np.zeros((N,N),dtype=complex)

    for j in range(N):
        for k in range(N):
            F[j,k]=omega**(j*k)

    return F/np.sqrt(N)
\end{lstlisting}

\end{frame}

%------------------------------------------------

\begin{frame}[fragile]{Example}

\begin{lstlisting}
n=3
F=qft_matrix(n)

state=np.zeros(2**n)
state[1]=1

result=F@state
print(result)
\end{lstlisting}

\end{frame}

%================================================
\section{Summary}

\begin{frame}{Summary}

\begin{itemize}
\item QFT efficiently computes Fourier transforms
\item central primitive in quantum algorithms
\item phase estimation extracts eigenvalues
\item widely used in quantum simulation
\end{itemize}

\end{frame}

\end{document}
